%============================================================
% CH1 · PATTERN 5 — FUNCTIONAL GRAPHS (ONE OUTGOING EDGE): THEORY
%============================================================
\subsection{Pattern Theory}
\label{ssec:ch1-p5-theory}

\begin{patternbox}{Pattern 5 at a Glance}
Every vertex has \textbf{exactly one outgoing edge} $v \to f(v)$, so the whole graph is just an array \texttt{f[]}. There is nothing to choose: from a start vertex the walk $v, f(v), f(f(v)), \dots$ is forced. Because there are only $n$ vertices the walk must repeat, and from then on it goes round and round. Every walk looks like the Greek letter $\rho$ (rho): a \textbf{tail} of $\mu$ vertices followed by a \textbf{cycle} of $\lambda$ vertices.

\begin{center}
\begin{tikzpicture}[x=1.1cm, y=1.1cm]
    \node[gvsrc] (s) at (0,0) {$s$};
    \node[gv] (t1) at (1.2,0) {};
    \node[gv] (t2) at (2.4,0) {};
    \node[gvdst] (e) at (3.6,0) {$e$};
    \node[gvseen] (c1) at (4.5,0.75) {};
    \node[gvseen] (c2) at (5.6,0.45) {};
    \node[gvseen] (c3) at (5.6,-0.45) {};
    \node[gvseen] (c4) at (4.5,-0.75) {};
    \draw[gde] (s) -- (t1); \draw[gde] (t1) -- (t2); \draw[gde] (t2) -- (e);
    \draw[gde] (e) -- (c1); \draw[gde] (c1) -- (c2); \draw[gde] (c2) -- (c3);
    \draw[gde] (c3) -- (c4); \draw[gde] (c4) -- (e);
    \draw[decorate, decoration={brace, amplitude=5pt, mirror}, darkgray] (-0.3,-0.55) -- (3.2,-0.55)
        node[midway, below=5pt, font=\scriptsize]{tail: $\mu = 3$ vertices};
    \node[font=\scriptsize] at (5.05,-1.25) {cycle: $\lambda = 5$};
    \node[note, text width=4.3cm, align=left, anchor=west] at (6.4,0)
        {$e$ = \textbf{entry point}: the first vertex the walk visits twice. The walk from $s$ visits $\mu+\lambda \le n$ different vertices.};
\end{tikzpicture}
\end{center}
\begin{recurrencebox}
\centering\small
walk from $s$ and remember the step $\texttt{first}[v]$ at which each vertex is reached; at the first vertex $v$ seen twice:\\[3pt]
$\mu = \texttt{first}[v]$, \qquad $\lambda = \text{step} - \texttt{first}[v]$,\\[2pt]
$f^{k}(s) = $ the vertex at step $\mu + (k-\mu) \bmod \lambda$ \;(for $k \ge \mu$).\\[3pt]
All vertices at once: walk from every new vertex with three colours \emph{new / on this walk / finished} --- \bigO{n} in total.
\end{recurrencebox}
\end{patternbox}

%------------------------------------------------------------
\subsubsection*{From Pattern 4 to Pattern 5}

In Pattern~4 a state had \emph{several} possible next states (press button A or button B, turn any of the four wheels), and we needed BFS to explore all of them. Now look at processes with \textbf{no choice at all}: a teacher who always goes to the student just accused, a teleporter that always sends you to the same town, a pseudo-random generator $x \mapsto (x^2+1) \bmod m$, a linked list whose \texttt{next} pointers may loop. Each state has exactly one successor --- the state graph has out-degree~1. Such a graph is called a \textbf{functional graph}, because its edges are the graph of a function $f : V \to V$ (Lesson~\ref{ssec:ch1-vocabulary}).

With no choices there is no search: we simply \emph{follow the arrows}. The interesting questions are about the \emph{shape} the arrows form. Where does the walk start repeating? How long is the loop? Which vertices lie on loops at all? Where will I be after $10^{18}$ steps?

\subsubsection*{From an array to a picture}

The input of a functional-graph problem is almost always a single array: \texttt{p[i]}, \texttt{next[i]}, \texttt{a[i]}, \texttt{favorite[i]}, \texttt{edges[i]}. Index $i$ is a vertex, the value is where its arrow points. No adjacency list is needed --- the array \emph{is} the adjacency list.

\begin{diagrambox}
\begin{tikzpicture}
    \dpindex{ix}{1.5}{$v$}{0}{9}
    \dparray{fa}{0.8}{$f[v]$}{1/cell, 2/cell, 3/cell, 4/cell, 2/cell, 4/cell, 5/cell, 8/cell, 9/cell, 8/cell}
    \begin{scope}[yshift=-1.9cm, xshift=3.6cm]
        \foreach \v/\x/\y in {0/-2.8/0.8, 1/-1.4/0.8, 2/0/0.8, 3/1.3/0, 4/0/-0.8, 5/-1.4/-0.8, 6/-2.8/-0.8, 7/2.9/0, 8/4.2/0, 9/5.5/0}
            \node[gv] (n\v) at (\x,\y) {\v};
        \foreach \u/\w in {0/1, 1/2, 2/3, 3/4, 4/2, 5/4, 6/5, 7/8}
            \draw[gde] (n\u) -- (n\w);
        \draw[gde] (n8) to[bend left=35] (n9);
        \draw[gde] (n9) to[bend left=35] (n8);
    \end{scope}
\end{tikzpicture}
\figcap{The array $f = [1,2,3,4,2,4,5,8,9,8]$ and its functional graph: 10 vertices, 10 arrows, two pieces. Each piece is one cycle ($2\to3\to4\to2$ and $8 \leftrightarrows 9$) with trees hanging into it.}
\end{diagrambox}

We call this example graph $G$ and use it throughout the section.

%------------------------------------------------------------
\subsubsection*{The shape theorem}

\begin{theorembox}[Shape of a functional graph]
\begin{enumerate}[leftmargin=*]
    \item \textbf{Every walk is a $\rho$.} From any start $s$ the walk $s, f(s), f^2(s), \dots$ visits $\mu$ tail vertices and then cycles forever through $\lambda$ cycle vertices, with $\mu \ge 0$, $\lambda \ge 1$ and $\mu + \lambda \le n$.
    \item \textbf{One cycle per piece.} Every connected piece (ignoring arrow directions) contains \emph{exactly one} cycle; all its other vertices form trees whose arrows point towards the cycle.
\end{enumerate}
\end{theorembox}

\textbf{Proof of 1.} Among the $n+1$ vertices $f^0(s), f^1(s), \dots, f^n(s)$ two must be equal (pigeonhole). Let $j$ be the first step that reaches an already visited vertex, say $f^j(s) = f^i(s)$ with $i < j$. The walk is deterministic, so from step $i$ on it repeats the steps $i, \dots, j-1$ forever. Set $\mu = i$ and $\lambda = j - i$; the vertices at steps $0..j-1$ are all different, so $\mu + \lambda = j \le n$.

\textbf{Proof of 2.} A piece with $k$ vertices has exactly $k$ edges (one leaves each vertex). A connected graph with $k$ vertices and $k-1$ edges is a tree; one more edge closes exactly one cycle (Lesson~\ref{ssec:ch1-vocabulary}). By part 1 the walk from any vertex of the piece ends in a \emph{directed} cycle, so that unique cycle is a directed one, and every vertex of the piece flows into it. \hfill$\blacksquare$

\begin{diagrambox}
\begin{tikzpicture}[x=1cm, y=1cm]
    % (a) permutation
    \begin{scope}
        \node[font=\small\bfseries] at (1.2,1.75) {(a) permutation};
        \foreach \v/\x/\y in {0/0/1, 1/1.1/1, 2/0.55/0.1, 3/1.9/0.1, 4/2.4/1}
            \node[gv] (a\v) at (\x,\y) {\v};
        \draw[gde] (a0) -- (a1); \draw[gde] (a1) -- (a2); \draw[gde] (a2) -- (a0);
        \draw[gde] (a3) to[bend left=30] (a4); \draw[gde] (a4) to[bend left=30] (a3);
        \node[note, text width=3.8cm] at (1.2,-0.85) {in-degree also 1:\\ \textbf{only cycles}, no tails};
    \end{scope}
    % (b) functional
    \begin{scope}[xshift=4.9cm]
        \node[font=\small\bfseries] at (1.2,1.75) {(b) functional graph};
        \foreach \v/\x/\y in {0/0/1, 1/1.1/1, 2/2.2/1, 3/2.2/0, 4/1.1/0}
            \node[gv] (b\v) at (\x,\y) {\v};
        \draw[gde] (b0) -- (b1); \draw[gde] (b1) -- (b2); \draw[gde] (b2) -- (b3);
        \draw[gde] (b3) -- (b4); \draw[gde] (b4) -- (b1);
        \node[note, text width=3.8cm] at (1.2,-0.85) {out-degree exactly 1:\\ \textbf{every piece is a $\rho$-shape} (one cycle + trees)};
    \end{scope}
    % (c) out-degree at most 1
    \begin{scope}[xshift=9.8cm]
        \node[font=\small\bfseries] at (1.2,1.75) {(c) out-degree $\le 1$};
        \foreach \v/\x/\y in {0/0/1, 1/1.1/1, 3/0/0, 4/1.1/0}
            \node[gv] (c\v) at (\x,\y) {\v};
        \node[gvdst] (c2) at (2.2,0.5) {2};
        \draw[gde] (c0) -- (c1); \draw[gde] (c1) -- (c2);
        \draw[gde] (c3) -- (c4); \draw[gde] (c4) -- (c2);
        \draw[-{Stealth[length=2mm]}, warningred, dashed] (c2) -- (3.0,0.5);
        \node[font=\scriptsize\bfseries, text=warningred, anchor=west] at (3.0,0.5) {$-1$};
        \node[note, text width=3.8cm] at (1.2,-0.85) {some \texttt{f[v] = -1}: a piece may be a \textbf{tree ending in a sink} (no cycle)};
    \end{scope}
\end{tikzpicture}
\figcap{Three relatives. Permutations are the special case where every vertex also has in-degree 1. In the ``at most one'' version (LeetCode writes \texttt{edges[i] = -1}) each piece has one cycle or none.}
\end{diagrambox}

%------------------------------------------------------------
\subsubsection*{When to recognise it (and when not)}

\begin{itemize}[leftmargin=*]
    \item the input is \textbf{one array} \texttt{a[1..n]} with ``$i$ points to / reports / likes / teleports to / is followed by $a_i$'';
    \item a \textbf{deterministic process}: the next state depends only on the current one ($x \mapsto g(x)$), and the question is about repeating, looping, or the state after $k$ steps;
    \item words like ``eventually'', ``forever'', ``the first student visited twice'', ``after $K = 10^{18}$ moves'', ``longest cycle'', ``everyone has exactly one favourite''.
\end{itemize}
\textbf{Not this pattern:} if a vertex can have \emph{two or more} outgoing edges, the walk is not forced; cycle detection then needs a real DFS with colours (Chapter~2 (DFS)) and the cycle structure is studied with strongly connected components later in the book. If the graph is undirected, ``one cycle per piece'' is a different object (a pseudotree, Chapter~13 (trees)). And if there are \emph{many} queries ``where am I after $k$ steps from $v$?'', walking once per query is too slow --- that is binary lifting, Chapter~11 (functional graphs and binary lifting); we only preview the idea below.

%------------------------------------------------------------
\subsubsection*{Algorithm 1 --- one start: measure the $\rho$}

Walk from $s$ and write down at which step each vertex is reached. The first time we arrive at a vertex that already has a step number, we have found the entry point $e$, the tail length $\mu = \texttt{first}[e]$ and the cycle length $\lambda = \text{step} - \texttt{first}[e]$.

\begin{diagrambox}
\begin{tikzpicture}[x=1.45cm]
    \foreach \v/\i/\s in {0/0/gvsrc, 1/1/gv, 2/2/gvdst, 3/3/gvseen, 4/4/gvseen}{
        \node[\s] (w\i) at (\i,0) {\v};
        \node[gdist, below=1mm of w\i] {step \i};
    }
    \node[gvdst, dashed] (w5) at (5,0) {2};
    \node[gdist, below=1mm of w5, text=warningred] {step 5};
    \foreach \i [evaluate=\i as \j using int(\i+1)] in {0,...,4} \draw[gde] (w\i) -- (w\j);
    \draw[gback] (w5.north) to[bend right=35] node[lbl, above]{seen before, at step 2} (w2.north);
    \draw[decorate, decoration={brace, amplitude=4pt, mirror}, darkgray] (-0.25,-0.85) -- (1.25,-0.85)
        node[midway, below=4pt, font=\scriptsize]{$\mu = \texttt{first}[2] = 2$};
    \draw[decorate, decoration={brace, amplitude=4pt, mirror}, darkgray] (1.75,-0.85) -- (4.25,-0.85)
        node[midway, below=4pt, font=\scriptsize]{$\lambda = 5 - 2 = 3$};
\end{tikzpicture}
\figcap{The walk from $0$ in $G$: $0,1,2,3,4$ and then $2$ again. Entry point $e = 2$, tail $\mu = 2$, cycle $\lambda = 3$. From $6$ the walk is $6,5,4,2,3$ and then $4$: entry point $4$, again $\mu = 2$, $\lambda = 3$ --- same cycle, different entry point.}
\end{diagrambox}

\textbf{Huge $k$.} Once we know $\mu$ and $\lambda$, the position after $k$ steps needs no long simulation: after the tail the walk is periodic with period $\lambda$, so for $k \ge \mu$ we may replace $k$ by $\mu + (k - \mu) \bmod \lambda$, which is smaller than $\mu + \lambda \le n$. In $G$ from $0$: $k = 10$ becomes $2 + 8 \bmod 3 = 4$, and step $4$ of the walk is vertex $4$; even $k = 10^{18}$ becomes $2 + (10^{18}-2) \bmod 3 = 4$, so $f^{10^{18}}(0) = 4$.

%------------------------------------------------------------
\subsubsection*{Algorithm 2 --- all vertices: walk with three colours}

Running Algorithm~1 from every vertex costs up to $\bigO{n}$ per start, $\bigO{n^2}$ in total --- fine for $n = 1000$, hopeless for $n = 2\cdot 10^5$. The fix is the idea behind \texttt{visited} in Lesson~\ref{ssec:ch1-first-traversal}: \textbf{never walk over a vertex twice}. Give each vertex one of three states:
\begin{center}
\legendchip{gv}{0 = new} \qquad \legendchip{gvseen}{1 = on the current walk} \qquad \legendchip{gvdone}{2 = finished}
\end{center}
For each start $s$ that is still new, walk while the vertices are new, colouring them 1 and storing them in \texttt{path}. The walk stops at the first vertex $v$ that is not new (or at $-1$):
\begin{itemize}[leftmargin=*]
    \item $v$ has \textbf{state 1}: we bit our own tail --- the part of \texttt{path} from $v$ to the end is a \textbf{new cycle};
    \item $v$ has \textbf{state 2}: we ran into territory explored by an earlier walk; its cycle was already found, so there is \textbf{no new cycle};
\end{itemize}
then colour the whole \texttt{path} 2.

\begin{diagrambox}
\begin{tikzpicture}
    % snapshot 1
    \begin{scope}
        \foreach \v/\x/\y/\s in {0/-2.8/0.8/gvseen, 1/-1.4/0.8/gvseen, 2/0/0.8/gvseen, 3/1.3/0/gvseen, 4/0/-0.8/gvseen,
                                 5/-1.4/-0.8/gv, 6/-2.8/-0.8/gv, 7/2.9/0/gv, 8/4.2/0/gv, 9/5.5/0/gv}
            \node[\s] (a\v) at (\x,\y) {\v};
        \foreach \u/\w in {0/1, 1/2, 2/3, 3/4, 5/4, 6/5, 7/8} \draw[gde] (a\u) -- (a\w);
        \draw[gback] (a4) -- (a2);
        \draw[gde] (a8) to[bend left=35] (a9); \draw[gde] (a9) to[bend left=35] (a8);
        \node[note, text width=5.2cm, align=left, anchor=west] at (6.3,0) {\textbf{$s=0$:} \texttt{path} $= 0,1,2,3,4$; the walk stops at $2$, which has state 1 $\Rightarrow$ new cycle $2,3,4$ of length $3$ (red arrow closes it).};
    \end{scope}
    % snapshot 2
    \begin{scope}[yshift=-2.75cm]
        \foreach \v/\x/\y/\s in {0/-2.8/0.8/gvdone, 1/-1.4/0.8/gvdone, 2/0/0.8/gvdone, 3/1.3/0/gvdone, 4/0/-0.8/gvdone,
                                 5/-1.4/-0.8/gvseen, 6/-2.8/-0.8/gv, 7/2.9/0/gv, 8/4.2/0/gv, 9/5.5/0/gv}
            \node[\s] (b\v) at (\x,\y) {\v};
        \foreach \u/\w in {0/1, 1/2, 2/3, 3/4, 4/2, 6/5, 7/8} \draw[gde] (b\u) -- (b\w);
        \draw[gcross] (b5) -- (b4);
        \draw[gde] (b8) to[bend left=35] (b9); \draw[gde] (b9) to[bend left=35] (b8);
        \node[note, text width=5.2cm, align=left, anchor=west] at (6.3,0) {\textbf{$s=5$:} \texttt{path} $= 5$; the walk stops at $4$, which has state 2 $\Rightarrow$ old territory, \textbf{no} new cycle. The same happens for $s = 6$ (stops at $5$).};
    \end{scope}
    % snapshot 3
    \begin{scope}[yshift=-5.5cm]
        \foreach \v/\x/\y/\s in {0/-2.8/0.8/gvdone, 1/-1.4/0.8/gvdone, 2/0/0.8/gvdone, 3/1.3/0/gvdone, 4/0/-0.8/gvdone,
                                 5/-1.4/-0.8/gvdone, 6/-2.8/-0.8/gvdone, 7/2.9/0/gvseen, 8/4.2/0/gvseen, 9/5.5/0/gvseen}
            \node[\s] (c\v) at (\x,\y) {\v};
        \foreach \u/\w in {0/1, 1/2, 2/3, 3/4, 4/2, 5/4, 6/5, 7/8} \draw[gde] (c\u) -- (c\w);
        \draw[gde] (c8) to[bend left=35] (c9);
        \draw[gback] (c9) to[bend left=35] (c8);
        \node[note, text width=5.2cm, align=left, anchor=west] at (6.3,0) {\textbf{$s=7$:} \texttt{path} $= 7,8,9$; the walk stops at $8$, state 1 $\Rightarrow$ new cycle $8,9$ of length $2$. Starts $8$ and $9$ are skipped later (already finished).};
    \end{scope}
\end{tikzpicture}
\figcap{Algorithm 2 on $G$. Orange = state 1 (current walk), blue = state 2 (finished), white = new. Each vertex is put on a \texttt{path} exactly once; result \texttt{cyc} $= [0,0,3,3,3,0,0,0,2,2]$.}
\end{diagrambox}

\begin{invariantbox}[Why Algorithm 2 is correct and linear]
\textbf{Invariant:} at every moment the vertices with state 1 are exactly the vertices of \texttt{path}, in walk order, and every cycle made of finished (state 2) vertices has already been reported.

\textbf{New cycle $\Leftrightarrow$ state 1.} If the walk stops at a state-1 vertex $v$, then $v$ is in \texttt{path}, and following arrows from $v$ along \texttt{path} leads back to $v$: that is a cycle, and it is new because its vertices were new when this walk started. If the walk stops at a state-2 vertex, every vertex reachable from it is finished too (an earlier walk followed exactly these arrows), so the $\rho$ of the current start ends in an old cycle; by the shape theorem there is no other cycle it could reach.

\textbf{Linear time:} a vertex is appended to \texttt{path} only while it is new, and then it leaves state 0 forever --- at most $n$ appends in total. Finding where the cycle starts scans the current \texttt{path} once. Total \bigO{n} time and memory.
\end{invariantbox}

%------------------------------------------------------------
\subsubsection*{Algorithm 3 --- peeling: who is on a cycle?}

A vertex with in-degree $0$ cannot lie on a cycle (a cycle would enter it). Delete it; this may lower the in-degree of its successor to $0$; delete that one too; and so on --- this is Kahn's queue from Chapter~5 (DAGs). What survives are precisely the cycle vertices.

\begin{diagrambox}
\begin{tikzpicture}
    \begin{scope}
        \foreach \v/\x/\y/\s in {0/-2.8/0.8/gvseen, 1/-1.4/0.8/gv, 2/0/0.8/gv, 3/1.3/0/gv, 4/0/-0.8/gv,
                                 5/-1.4/-0.8/gv, 6/-2.8/-0.8/gvseen, 7/2.9/0/gvseen, 8/4.2/0/gv, 9/5.5/0/gv}
            \node[\s] (a\v) at (\x,\y) {\v};
        \foreach \u/\w in {0/1, 1/2, 2/3, 3/4, 4/2, 5/4, 6/5, 7/8} \draw[gde] (a\u) -- (a\w);
        \draw[gde] (a8) to[bend left=35] (a9); \draw[gde] (a9) to[bend left=35] (a8);
        \node[note, text width=5.2cm, align=left, anchor=west] at (6.3,0) {in-degrees $[0,1,2,1,2,1,0,0,2,1]$;\\ the queue starts with $0, 6, 7$ (orange).};
    \end{scope}
    \begin{scope}[yshift=-2.75cm]
        \foreach \v/\x/\y/\s in {0/-2.8/0.8/gvdead, 1/-1.4/0.8/gvseen, 2/0/0.8/gv, 3/1.3/0/gv, 4/0/-0.8/gv,
                                 5/-1.4/-0.8/gvseen, 6/-2.8/-0.8/gvdead, 7/2.9/0/gvdead, 8/4.2/0/gv, 9/5.5/0/gv}
            \node[\s] (b\v) at (\x,\y) {\v};
        \foreach \u/\w in {1/2, 2/3, 3/4, 4/2, 5/4} \draw[gde] (b\u) -- (b\w);
        \foreach \u/\w in {0/1, 6/5, 7/8} \draw[gfaded, dashed] (b\u) -- (b\w);
        \draw[gde] (b8) to[bend left=35] (b9); \draw[gde] (b9) to[bend left=35] (b8);
        \node[note, text width=5.2cm, align=left, anchor=west] at (6.3,0) {peel $0, 6, 7$: the in-degrees of $1$ and $5$ drop to $0$, they join the queue; $8$ drops to $1$ only.};
    \end{scope}
    \begin{scope}[yshift=-5.5cm]
        \foreach \v/\x/\y/\s in {0/-2.8/0.8/gvdead, 1/-1.4/0.8/gvdead, 2/0/0.8/gvsel, 3/1.3/0/gvsel, 4/0/-0.8/gvsel,
                                 5/-1.4/-0.8/gvdead, 6/-2.8/-0.8/gvdead, 7/2.9/0/gvdead, 8/4.2/0/gvsel, 9/5.5/0/gvsel}
            \node[\s] (c\v) at (\x,\y) {\v};
        \foreach \u/\w in {2/3, 3/4, 4/2} \draw[gde] (c\u) -- (c\w);
        \foreach \u/\w in {0/1, 1/2, 5/4, 6/5, 7/8} \draw[gfaded, dashed] (c\u) -- (c\w);
        \draw[gde] (c8) to[bend left=35] (c9); \draw[gde] (c9) to[bend left=35] (c8);
        \node[note, text width=5.2cm, align=left, anchor=west] at (6.3,0) {peel $1, 5$: $2$ and $4$ keep in-degree $1$; the queue is empty. Survivors (purple) $2,3,4,8,9$ = the cycle vertices.};
    \end{scope}
\end{tikzpicture}
\figcap{Peeling $G$. Grey vertices are deleted, in the order $0, 6, 7, 1, 5$.}
\end{diagrambox}

\textbf{Why it works.} A cycle vertex always keeps the arrow from its cycle predecessor, which is never deleted (induction on the deletion order), so its in-degree never reaches 0. A tree vertex hangs in a tree whose arrows point to the cycle; the vertices farthest from the cycle have in-degree 0, and peeling them exposes the next layer --- every tree vertex is eventually deleted. Peeling is handy when information must \textbf{flow along the trees}, e.g.\ the longest chain that ends at each cycle vertex (Exercise E5.7), because a vertex is always deleted \emph{before} its successor.

%------------------------------------------------------------
\subsubsection*{Algorithm 4 --- Floyd's tortoise and hare (\bigO{1} memory)}

Two pointers start at $s$; per round the tortoise moves 1 step and the hare 2 steps. Inside the cycle the hare gains one step per round, so they meet. Then put the tortoise back at $s$ and move both one step at a time: they meet exactly at the entry point $e$, after $\mu$ steps. One more lap around the cycle gives $\lambda$.

\begin{diagrambox}
\begin{tikzpicture}
    \node[font=\small] at (0,0) {
    \begin{tabular}{@{}lcccc@{}}
    \toprule
    round $t$ & 0 & 1 & 2 & 3 \\
    \midrule
    tortoise $f^{t}(0)$ & 0 & 1 & 2 & \textbf{3} \\
    hare $f^{2t}(0)$ & 0 & 2 & 4 & \textbf{3} \\
    \bottomrule
    \end{tabular}};
    \node[note, text width=6.6cm, align=left, anchor=west] at (3.0,0)
        {In $G$ from $s = 0$ they meet at vertex $3$ after 3 rounds. Restart the tortoise at $0$, both move 1 step: $(0,3) \to (1,4) \to (2,2)$ --- they meet at the entry point $2$ after $\mu = 2$ steps; one lap from $2$ gives $\lambda = 3$.};
\end{tikzpicture}
\end{diagrambox}

Why the second phase works: when they first meet after $t$ rounds, the hare has walked $t$ more steps than the tortoise, and that surplus is a whole number of laps, a multiple of $\lambda$. So a pointer that walks $\mu$ more steps from the meeting point stands at step $t + \mu$ of the walk, which is on the cycle and equivalent to step $\mu$ modulo $\lambda$: the entry point --- exactly where the restarted tortoise arrives after $\mu$ steps. Floyd's method is the classic answer to ``detect a loop in a linked list with \bigO{1} memory''; in contests memory is rarely that tight, and the \texttt{first[]} array of Algorithm~1 is simpler. When the input is read-only and only \bigO{1} extra memory is allowed, Floyd is the tool; \hyperref[ssec:ch1-find-duplicate]{Problem~1.5.3} (Find the Duplicate Number) is such a case and proves the second phase in detail.

%------------------------------------------------------------
\subsubsection*{A preview: many $k$-th successor queries}

Algorithm~1 answers \emph{one} question ``where is $f^k(s)$?'' in \bigO{n}. With $q = 2\cdot10^5$ different starts that is too slow. The idea of Chapter~11 (functional graphs and binary lifting) is \textbf{doubling}: precompute $f^{1}, f^{2}, f^{4}, f^{8}, \dots$ for every vertex (using $f^{2^{j+1}}(v) = f^{2^j}(f^{2^j}(v))$), then write $k$ in binary and jump: $f^{13} = f^{8} \circ f^{4} \circ f^{1}$ --- only about $\log_2 k$ jumps per query.

\begin{diagrambox}
\begin{tikzpicture}[x=0.95cm]
    \foreach \i in {0,...,13}{
        \ifnum\i=0 \node[gvsrc, minimum size=6mm, font=\scriptsize\bfseries] (q\i) at (\i,0) {\i};
        \else\ifnum\i=13 \node[gvdst, minimum size=6mm, font=\scriptsize\bfseries] (q\i) at (\i,0) {\i};
        \else \node[gv, minimum size=6mm, font=\scriptsize\bfseries] (q\i) at (\i,0) {\i}; \fi\fi
    }
    \foreach \i [evaluate=\i as \j using int(\i+1)] in {0,...,12} \draw[gfaded, -{Stealth[length=1.6mm]}] (q\i) -- (q\j);
    \draw[gpath, -{Stealth[length=3mm]}] (q0.north) to[bend left=30] node[lbl, above]{jump $2^3 = 8$} (q8.north);
    \draw[gpath, -{Stealth[length=3mm]}] (q8.north) to[bend left=40] node[lbl, above]{$2^2 = 4$} (q12.north);
    \draw[gpath, -{Stealth[length=3mm]}] (q12.south) to[bend right=70] node[lbl, below]{$2^0 = 1$} (q13.south);
\end{tikzpicture}
\figcap{Doubling: 13 single steps (grey) become three precomputed jumps, $13 = 8 + 4 + 1$. The walk is drawn as a line of step numbers; in a functional graph the jumps simply wrap around the cycle. Details and proofs: Chapter~11.}
\end{diagrambox}

%------------------------------------------------------------
\subsubsection*{Variations map}

\begin{diagrambox}
\begin{tikzpicture}[
    var/.style={rectangle, rounded corners=2pt, draw=patternborder, fill=patternbg,
                font=\scriptsize, align=center, text width=3.25cm, minimum height=1.3cm, inner sep=3pt}]
    \node[var, fill=answerfill, draw=accentorange, text width=5cm, font=\scriptsize\bfseries] (root) at (0,0)
        {Follow the arrows $v \to f(v)$\\every piece = one cycle + trees};
    \node[var] (v1) at (-5.4,-2) {\textbf{first repeated vertex}\\from every start\\(Badge, 1.5.1)};
    \node[var] (v2) at (-1.8,-2) {\textbf{longest cycle}\\out-degree $\le 1$\\(Longest Cycle, 1.5.2)};
    \node[var] (v3) at (1.8,-2)  {\textbf{permutation}\\only cycles: lengths, lcm\\(E5.1, E5.5)};
    \node[var] (v4) at (5.4,-2)  {\textbf{position after $K$ steps}\\$K$ huge: skip laps\\(E5.4)};
    \node[var] (v5) at (-5.4,-3.6) {\textbf{implicit $f$}\\$x \mapsto g(x)$: $\rho$ of a generator\\(E5.2)};
    \node[var] (v6) at (-1.8,-3.6) {\textbf{distances along walks}\\where two walkers meet\\(E5.3)};
    \node[var] (v7) at (1.8,-3.6)  {\textbf{one choice per cycle}\\cheapest vertex of each cycle\\(E5.6)};
    \node[var] (v8) at (5.4,-3.6)  {\textbf{cycles + longest tails}\\peeling, chain lengths, layers\\(E5.7, E5.8)};
    \node[var, text width=5cm] (v9) at (0,-5.2) {\textbf{entry point with \bigO{1} memory}\\read-only array: Floyd's two pointers\\(Find the Duplicate, 1.5.3)};
    \node[draw=patternborder, dashed, rounded corners=4pt, fit=(v1)(v4)(v5)(v8)(v9), inner sep=5pt] (grp) {};
    \draw[flow] (root.south) -- (grp.north);
    \node[note, below=2mm of grp] {many $k$-th successor queries (``after $k$ steps'' for many starts) $\Rightarrow$ binary lifting, Chapter~11};
\end{tikzpicture}
\end{diagrambox}

\textbf{Twists and how to handle them}
\begin{itemize}[leftmargin=*]
    \item \textbf{Missing arrows} (\texttt{f[v] = -1}): stop the walk at $-1$; such a walk contains no cycle.
    \item \textbf{Self-loops} (\texttt{f[v] = v}) are cycles of length 1: a student who blames himself gets his second hole immediately.
    \item \textbf{1-indexed input}: allocate $n+1$ entries, or subtract 1 while reading --- but do it consistently.
    \item \textbf{``Something per cycle''} (cheapest vertex, cycle length, number of cycles): Algorithm 2 reports each cycle exactly once.
    \item \textbf{``Something per vertex''} (its entry point, its distance to the cycle): first find the cycles, then fill the tails from the cycle outwards (Problem~1.5.1).
    \item \textbf{Read-only input, \bigO{1} extra memory}: no \texttt{first[]} or \texttt{state[]} array is allowed; walk with two pointers (Algorithm 4, \hyperref[ssec:ch1-find-duplicate]{Problem~1.5.3}).
\end{itemize}

%------------------------------------------------------------
\subsubsection*{Code templates}

\begin{lstlisting}[style=cpp]
#include <bits/stdc++.h>
using namespace std;
typedef long long ll;

// (1) Walk from ONE start s: tail length mu and cycle length lam.
//     Needs out-degree exactly 1 (f[v] is always a vertex).
void rho(const vector<int>& f, int s, int& mu, int& lam) {
    vector<int> first(f.size(), -1);      // first[v] = step at which v was first reached
    int v = s, step = 0;
    while (first[v] == -1) {              // keep walking until a vertex repeats
        first[v] = step++;
        v = f[v];
    }
    mu = first[v];                        // v is the entry point of the cycle
    lam = step - first[v];                // steps between the two visits of v
}

// (2) k-th successor f^k(s) for huge k: skip whole laps of the cycle
int kth(const vector<int>& f, int s, ll k) {
    int mu, lam;
    rho(f, s, mu, lam);
    if (k > mu) k = mu + (k - mu) % lam;  // after the tail the walk has period lam
    int v = s;
    for (ll i = 0; i < k; i++) v = f[v];  // now k < mu + lam <= n
    return v;
}

// (3) ALL vertices: walk with three states. f[v] == -1 means "no outgoing edge".
//     cyc[v] = length of the cycle through v, or 0 if v is not on a cycle.
vector<int> cycleLengths(const vector<int>& f) {
    int n = f.size();
    vector<int> state(n, 0), cyc(n, 0);   // 0 = new, 1 = on the current walk, 2 = finished
    for (int s = 0; s < n; s++) {
        if (state[s] != 0) continue;
        vector<int> path;                 // the current walk, in order
        int v = s;
        while (v != -1 && state[v] == 0) {
            state[v] = 1;
            path.push_back(v);
            v = f[v];
        }
        if (v != -1 && state[v] == 1) {   // we bit our own tail: a NEW cycle
            int len = path.end() - find(path.begin(), path.end(), v);
            for (int i = (int)path.size() - len; i < (int)path.size(); i++) cyc[path[i]] = len;
        }
        for (int u : path) state[u] = 2;  // the whole walk is finished
    }
    return cyc;
}

// (4) Peeling: delete vertices of in-degree 0 again and again (Kahn).
//     The survivors are exactly the vertices that lie on cycles.
vector<bool> onCycle(const vector<int>& f) {
    int n = f.size();
    vector<int> indeg(n, 0);
    for (int v = 0; v < n; v++)
        if (f[v] != -1) indeg[f[v]]++;
    queue<int> q;
    for (int v = 0; v < n; v++)
        if (indeg[v] == 0) q.push(v);
    vector<bool> alive(n, true);
    while (!q.empty()) {
        int v = q.front(); q.pop();
        alive[v] = false;                 // nobody points to v: v is not on a cycle
        if (f[v] != -1 && --indeg[f[v]] == 0) q.push(f[v]);
    }
    return alive;
}

// (5) Floyd's tortoise and hare: mu and lam with O(1) extra memory
void floyd(const vector<int>& f, int s, int& mu, int& lam) {
    int t = f[s], h = f[f[s]];            // tortoise 1 step, hare 2 steps
    while (t != h) { t = f[t]; h = f[f[h]]; }
    mu = 0; t = s;                        // restart the tortoise; both now move 1 step
    while (t != h) { t = f[t]; h = f[h]; mu++; }
    lam = 1; h = f[t];                    // walk once around the cycle
    while (t != h) { h = f[h]; lam++; }
}

int main() {
    vector<int> f = {1, 2, 3, 4, 2, 4, 5, 8, 9, 8};
    int mu, lam;
    rho(f, 0, mu, lam);   cout << mu << " " << lam << "\n";            // 2 3
    rho(f, 6, mu, lam);   cout << mu << " " << lam << "\n";            // 2 3
    floyd(f, 7, mu, lam); cout << mu << " " << lam << "\n";            // 1 2
    cout << kth(f, 0, 10) << " " << kth(f, 0, 1000000000000000000LL) << "\n";  // 4 4
    vector<int> cyc = cycleLengths(f);
    for (int x : cyc) cout << x << " ";                                // 0 0 3 3 3 0 0 0 2 2
    cout << "\n";
    vector<bool> alive = onCycle(f);
    for (int v = 0; v < 10; v++) if (alive[v]) cout << v << " ";       // 2 3 4 8 9
    cout << "\n";
}
\end{lstlisting}

\begin{insightbox}[Which template when?]
\renewcommand{\arraystretch}{1.25}
\begin{tabularx}{\linewidth}{@{}>{\bfseries}l X@{}}
(1) \texttt{rho} & one start: where does it start repeating, how long is the loop? \\
(2) \texttt{kth} & one start: position after a huge number of steps \\
(3) \texttt{cycleLengths} & every cycle exactly once, or a value per vertex; works with \texttt{-1} \\
(4) \texttt{onCycle} & which vertices are on cycles, plus DP along the trees towards the cycles \\
(5) \texttt{floyd} & only \bigO{1} memory allowed (linked lists, generators you cannot store, read-only arrays as in \hyperref[ssec:ch1-find-duplicate]{Problem~1.5.3}) \\
\end{tabularx}
\end{insightbox}

%------------------------------------------------------------
\subsubsection*{Common mistakes}

\begin{warningbox}[Mistake 1 --- a fresh \texttt{visited} array for every start]
\begin{lstlisting}[style=snippet]
for (int s = 0; s < n; s++) {
    vector<bool> seen(n, false);          // O(n) to create, O(n) to walk: O(n^2) overall
    ...
}
\end{lstlisting}
Fine for $n \le 1000$ (Badge), \tle{} for $n = 10^5$. Keep \emph{one} state array for all walks (Algorithm 2).
\end{warningbox}

\begin{warningbox}[Mistake 2 --- ``I met a visited vertex, so I found a new cycle'']
With a two-state \texttt{visited} array, the walk from $5$ in $G$ stops at $4$ (visited by the walk from $0$) and would wrongly report a ``cycle'' through $5$. Only a vertex visited \emph{during the current walk} closes a new cycle: that is why we need three states (or visit times compared with the start time of the walk, Problem~1.5.2).
\end{warningbox}

\begin{warningbox}[{Mistake 3 --- forgetting \texttt{f[v] = -1}}]
In out-degree-$\le 1$ problems, \texttt{state[f[v]]} with \texttt{f[v] = -1} reads outside the array. Test \texttt{v != -1} \emph{before} every access.
\end{warningbox}

\begin{warningbox}[Mistake 4 --- simulating $K$ steps, or overflowing $K$]
$K = 10^{18}$ steps cannot be simulated. Reduce $K$ with $\mu$ and $\lambda$ first, and keep $K$ in \texttt{long long} --- including the expression $(K - \mu) \bmod \lambda$.
\end{warningbox}

\begin{warningbox}[Mistake 5 --- recursion along the arrows]
A recursive \texttt{solve(v)} that calls \texttt{solve(f[v])} recurses as deep as the longest tail plus cycle --- up to $n = 2\cdot10^5$ levels on one long chain, a classic \stackoverflow. Use a loop with an explicit \texttt{path} vector, as all templates above do.
\end{warningbox}

\begin{checklistbox}[``I'm stuck'' checklist for Pattern 5]
\begin{itemize}[leftmargin=*]
    \item[$\square$] Does every vertex have exactly one outgoing arrow (or at most one)? Then draw a few $\rho$'s.
    \item[$\square$] Is the question about \emph{one} start (Algorithm 1) or about \emph{all} vertices / all cycles (Algorithm 2 or 3)?
    \item[$\square$] Am I visiting each vertex once overall, not once per start?
    \item[$\square$] When the walk stops: at a vertex of \emph{this} walk (new cycle), of an older walk (old cycle) or at $-1$ (no cycle)?
    \item[$\square$] Is the answer a property of the cycle (length, minimum) or of the tail (entry point, distance)? Fill tails from the cycle outwards.
    \item[$\square$] Huge $K$? Reduce it modulo the cycle length after the tail. Many queries? Chapter~11.
\end{itemize}
\end{checklistbox}
